\section{6. Evolución prevista de los sistemas de Health y análisis de señales}
Como complemento de los sistemas de Health ya implementados para la infraestructura, los protocolos, las sesiones y los Readers, se tiene previsto ampliar el análisis hacia el contenido audiovisual de las señales entrantes. El objetivo es que EBP pueda evaluar no solamente si existe una conexión o si circulan datos, sino también si la señal recibida conserva las condiciones audiovisuales necesarias para su operación.

\begin{tabularx}{\textwidth}{|L{0.34\textwidth}|Y|}\hline
\rowcolor{Navy}\THead{Análisis previsto} & \THead{Condición que se busca detectar} \\ \hline
Presencia de video & Determinar si la señal entrante contiene video utilizable o si el componente de video ha desaparecido. \\ \hline
\rowgray Presencia de audio & Determinar si existe audio y detectar una eventual pérdida de este componente. \\ \hline
Congelamiento de imagen & Identificar una señal que continúa conectada pero cuya imagen permanece congelada. \\ \hline
\rowgray Pantalla negra & Detectar condiciones en las que el flujo continúa presente, pero el contenido de video se encuentra en negro. \\ \hline
Continuidad audiovisual & Relacionar el estado del transporte con la continuidad efectiva del audio y del video. \\ \hline
\rowgray Condiciones de red & Incorporar evidencia de red para reconocer degradaciones que puedan afectar la recepción o distribución de la señal. \\ \hline
Localización de la incidencia & Correlacionar la evidencia disponible para ayudar a determinar si una falla se origina en la infraestructura de EBP, en la red, en la fuente o en el extremo del cliente. \\ \hline
\end{tabularx}

Esta evolución busca reducir la diferencia entre una conexión técnicamente activa y una señal realmente saludable. Por ejemplo, una sesión podría continuar conectada y transmitiendo datos mientras el video está congelado, el audio ha desaparecido o la imagen permanece en negro. El análisis de contenido permitiría que estas condiciones formen parte de la supervisión operacional.

La identificación del origen de una incidencia se plantea mediante correlación de evidencia y no únicamente mediante el estado de una conexión. Esto permitirá disponer de más elementos para distinguir problemas internos de EBP, condiciones de red y situaciones asociadas a la fuente o al receptor del cliente.

\begin{callout}La evolución del Health apunta a responder dos preguntas operacionales: ¿la señal está técnicamente conectada? y ¿el contenido de audio y video está realmente en condiciones correctas?\end{callout}
